\section{Generación de \emph{traces} para entrenamiento}
\label{sec:traces}

Para el entrenamiento posterior del modelo objetivo (Qwen3.5-2B) se generaron
\emph{traces} con Qwen3.5-9B como \emph{teacher} sobre la partición de
entrenamiento. Esta sección describe cómo se generaron y qué datos resultaron;
el entrenamiento mismo no se ha ejecutado. Los datos se publican con el
repositorio en \texttt{data/sudoku/} (particiones \texttt{train}, \texttt{eval} y
\texttt{test}) y \texttt{data/traces/} (\texttt{sft.jsonl} y \texttt{dpo.jsonl}).

\subsection{Configuración de la generación}
\label{sec:traces-config}

\begin{table}[htbp]
  \centering\small
  \caption{Configuración de la generación de \emph{traces}.}
  \label{tab:traces-config}
  \begin{tabular}{@{}ll@{}}
    \toprule
    Parámetro & Valor \\
    \midrule
    Modelo \emph{teacher} & Qwen3.5-9B, sin modo \emph{thinking} \\
    Tareas & \texttt{train}: 4x4 (fácil, medio, difícil) y 6x6 (fácil, medio) \\
    & 50 tareas por (tamaño, nivel): 250 tareas \\
    Excluidas & 6x6 difícil y 9x9 \\
    Muestras por tarea & $n = 8$ (2.000 muestras en la ronda r0) \\
    Muestreo (r0) & temperatura 0,7, \texttt{top\_p} 0,8 \\
    Límite de salida & 4.096 tokens \\
    \emph{Few-shot} & 2 ejemplos, formato \texttt{steps} (\texttt{Step k: rXcY = v} + rejilla) \\
    \bottomrule
  \end{tabular}
\end{table}

\paragraph{Qué se excluyó y por qué.} La partición de entrenamiento contiene 450
tareas (50 por tamaño y nivel, incluidos 6x6 difícil y 9x9). Se usó solo el
subconjunto que el 9B puede resolver con un presupuesto de 4.096 tokens;
6x6 difícil y 9x9 quedaron fuera. Todas las tareas tienen solución única
demostrada y se resuelven por casillas con un único candidato, de modo que el
\emph{reward} del verificador es exacto.

\paragraph{Decisión de usar \emph{few-shot}.} Un piloto de 5 tareas por nivel
con $n = 8$ comparó 0 y 2 \emph{shots}. Con 2 \emph{shots} las respuestas son más
cortas y truncan menos (6x6 medio: 40\,\% frente a 52\,\% de truncadas, y
\texttt{pass@8} de 1,00 frente a 0,60), y siguen un formato uniforme
(\texttt{Step k: rXcY = v}, rejilla y \texttt{<answer>}) que un modelo de 2B puede
aprender. Cada celda del piloto tiene solo 5 tareas (40 muestras), por lo que el
resultado es indicativo. Los \emph{shots} se usan solo para generar: el
\emph{prompt} guardado en los datos de entrenamiento es el \emph{prompt} base de
\texttt{train.jsonl}, sin ejemplos, que es el que el 2B recibe en la evaluación.
Así el 2B aprende el estilo de los ejemplos sin necesitarlos en contexto.

\paragraph{Ejecución.} Cada ronda se repartió en 5 \emph{shards} (un trabajo de
Slurm por \emph{shard}, una GPU L40S cada uno), y todos reciben las mismas 10
tareas por (tamaño, nivel). La corrida r0 tardó entre 1 h 48 min y 1 h 57 min por
\emph{shard}. Una segunda ronda (r1) con temperatura 1,0 y otra semilla se ejecutó
solo sobre las 100 tareas de 6x6 (800 muestras, 1 h 5 min a 1 h 15 min por
\emph{shard}) para conseguir más rechazados con rejilla errónea.

\subsection{Resultados de la generación}
\label{sec:traces-resultados}

\begin{table}[htbp]
  \centering\small
  \caption{Ronda r0: 250 tareas, 2.000 muestras, Qwen3.5-9B con 2 \emph{shots} y
  límite de 4.096 tokens.}
  \label{tab:traces-r0}
  \begin{tabular}{@{}lcccc@{}}
    \toprule
    Tablero / nivel & \passuno & \texttt{pass@8} & Truncadas & Tokens medios \\
    \midrule
    4x4 fácil  & 0,99 & 1,00 &  0\,\% &  771 \\
    4x4 medio  & 0,98 & 1,00 &  1\,\% & 1.113 \\
    4x4 difícil& 0,96 & 1,00 &  4\,\% & 1.772 \\
    6x6 fácil  & 0,87 & 1,00 & 10\,\% & 2.127 \\
    6x6 medio  & 0,64 & 0,96 & 34\,\% & 3.053 \\
    \bottomrule
  \end{tabular}
\end{table}

De las 2.000 muestras, 1.775 son correctas (\emph{reward} 1,0), y 248 de las 250
tareas tienen al menos una; las dos sin ninguna son de 6x6 medio. Los fallos son
sobre todo truncamientos: de 225 muestras fallidas, 194 (86\,\%) se cortaron en
4.096 tokens y solo 31 son rejillas erróneas o formato roto. Por eso, los
rechazados reales de DPO enseñan sobre todo a ser más breve. Además, las trazas
correctas se concentran en lo fácil (cerca de 1.170 son 4x4 y 600 son 6x6). La
ronda r1 (temperatura 1,0) no cambió las tasas de aprobación ni de truncamiento
(6x6 fácil 0,91 y 9\,\% truncadas; 6x6 medio 0,63 y 36\,\%) y aportó solo 6
rechazados con rejilla errónea, por lo que no se hicieron más rondas con el 9B.

\subsection{Construcción de los conjuntos SFT y DPO}
\label{sec:traces-datasets}

El script \texttt{build\_sft\_dpo.py} lee los \emph{shards} de r0 y r1, recupera el
\emph{prompt} sin \emph{shots} de \texttt{train.jsonl} por \texttt{task\_id} y
rechaza muestras que no sean de la partición \texttt{train}.

\paragraph{SFT.} Cada ejemplo es un par (\texttt{prompt}, \texttt{completion}).
Se incluyen las muestras con \emph{reward} 1,0 y sin truncar, hasta 2
\emph{completions} distintas por tarea, las más cortas primero. El resultado son
498 ejemplos: 100 por nivel, salvo 98 en 6x6 medio (cubre 250 tareas distintas
menos las dos sin muestras correctas).

\paragraph{DPO.} Cada ejemplo tiene \texttt{prompt}, \texttt{chosen},
\texttt{rejected} y los indicadores \texttt{rejected\_truncated} y
\texttt{rejected\_synthetic}. El \texttt{chosen} es la correcta más corta, y hasta
2 pares por tarea sin repetir \emph{completions}. El \texttt{rejected} real
prefiere rejillas erróneas que terminaron (la más cercana a la correcta primero) y
solo después las truncadas.

\paragraph{Pares sintéticos.} Como el 9B casi no produce rejillas erróneas, se
añadió un par sintético por tarea: el \texttt{rejected} es la propia traza correcta
con una celda originalmente vacía cambiada en la rejilla del \texttt{<answer>}
final. Cada uno se comprueba con el \emph{reward} (formato correcto y no
resuelto) y la semilla depende de \texttt{task\_id}, así que el resultado es
determinista. Son rechazados fuera de política: el razonamiento previo queda
intacto y solo la rejilla final es incorrecta, de modo que el par enseña sobre todo
a copiar fielmente la última rejilla verificada. Los indicadores permiten
filtrarlos o ponderarlos.

\begin{table}[htbp]
  \centering\small
  \caption{Pares DPO por nivel (432 en total).}
  \label{tab:traces-dpo}
  \begin{tabular}{@{}lcccccc@{}}
    \toprule
    Origen & 4x4 fácil & 4x4 medio & 4x4 difícil & 6x6 fácil & 6x6 medio & Total \\
    \midrule
    Reales      &  3 &  6 & 16 & 65 & 92 & 182 \\
    Sintéticos  & 50 & 50 & 50 & 50 & 50 & 250 \\
    \bottomrule
  \end{tabular}
\end{table}

De los 182 pares reales, 36 tienen un rechazado con rejilla errónea y 146 uno
truncado. Solo con r0 y sin sintéticos habría 148 pares reales (30 con rejilla
errónea).

\subsection{Limitaciones}
\label{sec:traces-limitaciones}
Los datos están sesgados a lo fácil y el 86\,\% de los fallos del 9B son
truncamientos, no errores de razonamiento. Los pares sintéticos no son
\emph{on-policy}. Un paso opcional es muestrear el 2B tras el SFT y emparejar sus
fallos con las trazas correctas del 9B. Además, ningún \emph{prompt} de
\texttt{train.jsonl} debería aparecer en \texttt{eval.jsonl} ni en
\texttt{test.jsonl}: las semillas impiden compartir puzzles, pero esto debe
verificarse. Las cifras de esta sección provienen de los registros de la
generación y no de las evaluaciones de las secciones anteriores, que usan la
partición de prueba.
